\lysh{22261_SOLUTION}{1}{
ΛΥΣΗ

α) Από τα δεδομένα έχουμε: $ΑΒ= γ$, $ΑΓ= β$ και $ΒΓ= \sqrt{β^2 + γ^2 + βγ}$.
Είναι: $ΒΓ^2 = β^2 + γ^2 + βγ$ και $ΑΓ^2 + ΑΒ^2 = β^2 + γ^2$.
Οπότε: $ΒΓ^2 > ΑΓ^2 + ΑΒ^2 \Rightarrow \hat{A} > 90^\circ$, άρα το τρίγωνο ΑΒΓ είναι αμβλυγώνιο.

β) Στο τρίγωνο ΑΒΓ, από το νόμο των συνημιτόνων έχουμε:
\[ ΒΓ^2 = ΑΓ^2 + ΑΒ^2 - 2ΑΓ \cdot ΑΒ \cdot \sigma\upsilon\nuΑ = β^2 + γ^2 - 2βγ \cdot \sigma\upsilon\nuΑ . \]
Επίσης από το ερώτημα (α): $ΒΓ^2 = β^2 + γ^2 + βγ$.
Άρα: $β^2 + γ^2 - 2βγ \cdot \sigma\upsilon\nuΑ = β^2 + γ^2 + βγ$ ή $\sigma\upsilon\nuΑ = -\frac{1}{2}$, άρα $\hat{A} = 120^\circ$.

\begin{center}
\begin{tikzpicture}[scale=1]
    % Ορισμός σημείων
    \tkzDefPoint(0,0){O}
    \tkzDefPoint(150:2.5){B}
    \tkzDefPoint(30:2.5){G}
    \tkzDefPoint(90:2.5){A}
    \tkzDefPoint(-90:2.5){M}

    % Σχεδίαση κύκλου
    \tkzDrawCircle(O,A)
    
    % Σχεδίαση τριγώνου ΑΒΓ
    \tkzDrawPolygon(A,B,G)
    \tkzDrawSegment(B,G) % Η χορδή ΒΓ
    
    % Σημείωση γωνίας Α
    \tkzMarkAngle[size=0.6](G,A,B)
    \tkzLabelAngle[pos=0.4](G,A,B){$120^\circ$}
    
    % Σχεδίαση σημείων και ετικετών
    \tkzDrawPoints[size=3,fill=black](A,B,G,M,O)
    \tkzLabelPoint[above](A){A}
    \tkzLabelPoint[left](B){B}
    \tkzLabelPoint[right](G){Γ}
    \tkzLabelPoint[below](M){M}
    \tkzLabelPoint[below right=-1pt and -1pt](O){O}
\end{tikzpicture}
\end{center}

γ) Η γωνία ΒΑΓ είναι εγγεγραμμένη στον κύκλο, οπότε το τόξο στο οποίο βαίνει θα έχει μέτρο διπλάσιο από αυτήν. Άρα $\stackrel{\frown}{ΒΜΓ} = 2 \cdot 120^\circ = 240^\circ$.

Έτσι για το κυρτογώνιο τόξο ΒΑΓ θα ισχύει: $\stackrel{\frown}{ΒΑΓ} = 360^\circ - 240^\circ = 120^\circ$, οπότε η επίκεντρη γωνία ΒΟΓ ισούται με $120^\circ$.

δ) Λόγω του ερωτήματος (γ) είναι: $\widehat{ΒΟΓ} = 120^\circ$, οπότε το ζητούμενο εμβαδό είναι:
\[ E = (Ο \stackrel{\frown}{ΒΑΓ}) - (ΟΒΓ) = \frac{πR^2 \mu}{360} - \frac{1}{2}ΟΒ \cdot ΟΓ \cdot \eta\mu 120^\circ = \frac{πR^2 120}{360} - \frac{1}{2}R \cdot R \cdot \frac{\sqrt{3}}{2} = \]
\[ = \frac{πR^2}{3} - \frac{R^2\sqrt{3}}{4} = \frac{4πR^2 - 3\sqrt{3}R^2}{12} = \frac{(4π-3\sqrt{3})R^2}{12} . \]
}